\esSezione{Sistema PO/PC}\label{sec:esercizio6-1}

\begin{traccia}
Progettare, implementare in VHDL e verificare mediante simulazione un sistema composto da una memoria ROM di $N$ locazioni da 8 bit, una macchina combinatoria $M$ che trasforma la stringa letta dalla ROM in una stringa di 4 bit e una memoria MEM di $N$ locazioni che memorizza l'uscita di $M$. Il sistema parte a fronte di un segnale \texttt{START} esterno; una unità di controllo gestisce la tempificazione, scandisce una locazione della ROM alla volta e scrive la locazione corrispondente di MEM. Gli indirizzi sono forniti da un contatore; ROM e MEM hanno rispettivamente lettura e scrittura sincrone.
\end{traccia}

Rispetto alla soluzione proposta, l'implementazione si discosta in vari punti, indicati nel seguito: i nomi di entità e porte, la parametrizzazione con due generic (\texttt{N} e \texttt{A}), la MEM dotata di una porta di lettura per il testbench, l'unità di controllo con uscita \texttt{done} e senza segnale di reset del contatore, e la funzione scelta per $M$.

\progetto

Il sistema è diviso in parte operativa e parte di controllo. I componenti istanziati nel top module \texttt{sistema} sono:
\begin{enumerate}
  \item Control Unit (\texttt{control\_unit});
  \item ROM (\texttt{rom});
  \item macchina combinatoria $M$ (\texttt{onecount});
  \item memoria (\texttt{mem});
  \item contatore (\texttt{counter\_mod}).
\end{enumerate}

Il contatore fornisce l'indirizzo, comune a ROM e MEM. La ROM restituisce il byte letto su \texttt{rom\_data}; la macchina $M$ lo trasforma in \texttt{m\_data}, che è il dato in ingresso alla MEM. L'unità di controllo riceve \texttt{start} e il riporto \texttt{tc} del contatore e produce i comandi \texttt{rd} (lettura della ROM), \texttt{we} (scrittura della MEM), \texttt{cnt\_en} (avanzamento del contatore) e l'uscita \texttt{done}. La funzione scelta per $M$ è il conteggio degli '1' presenti nel byte, che dà un risultato tra 0 e 8, rappresentabile su 4 bit.

% TODO: schema a blocchi dell'architettura PO/PC non presente tra le immagini del progetto Vivado
Lo schema a blocchi dell'architettura non è ancora stato inserito.

\implementazione

Si parte dai componenti della parte operativa. La ROM è descritta dall'entità \texttt{rom}, con generic \texttt{N} (locazioni) e \texttt{A} (bit di indirizzo) e porte \texttt{clk}, \texttt{rd}, \texttt{addr} (\texttt{A} bit) e \texttt{data} (8 bit). L'architettura \texttt{behavioral} contiene la costante \texttt{CONTENT}, un array di \texttt{N} byte con sedici valori iniziali e gli altri a \texttt{x"00"}. Un process sensibile a \texttt{clk} carica \texttt{data} sul fronte di salita solo se \texttt{rd} vale '1', quindi la lettura è sincrona; un \texttt{assert} con severità \texttt{failure} verifica che \texttt{N} sia indirizzabile su \texttt{A} bit. Rispetto alla soluzione proposta, il contenuto della ROM è diverso e non compare l'attributo \texttt{rom\_style}.

\vhdlfile{esercizio6-1/rom.vhd}{Implementazione della ROM}{lst:esercizio6-1:rom}

La macchina combinatoria $M$ è l'entità \texttt{onecount}, con ingresso \texttt{x} a 8 bit e uscita \texttt{y} a 4 bit. Un process sensibile a \texttt{x} azzera una variabile \texttt{n} di tipo \texttt{unsigned} e la incrementa per ogni bit di \texttt{x} pari a '1', con un ciclo \texttt{for} sugli otto indici; \texttt{y} è il risultato convertito in \texttt{std\_logic\_vector}. Non essendoci il clock, la rete è puramente combinatoria.

\vhdlfile{esercizio6-1/onecount.vhd}{Implementazione della macchina combinatoria $M$ (conteggio degli '1')}{lst:esercizio6-1:onecount}

L'entità \texttt{mem} ha generic \texttt{N}, \texttt{A} e \texttt{W} (larghezza del dato, 4 per default) e porte \texttt{clk}, \texttt{we}, \texttt{waddr}, \texttt{din}, oltre a \texttt{raddr} e \texttt{dout}, aggiunte per consentire al testbench di leggere il contenuto. La scrittura avviene in un process sincrono, quando \texttt{we} vale '1'; la lettura su \texttt{dout} è invece un assegnamento concorrente, quindi asincrona. Le locazioni sono inizializzate a tutti '1'. A differenza della soluzione proposta, la MEM espone quindi un'uscita.

\vhdlfile{esercizio6-1/mem.vhd}{Implementazione della memoria MEM}{lst:esercizio6-1:mem}

Il contatore è \texttt{counter\_mod}, già presentato nell'esercizio 5.1 (listato~\ref{lst:esercizio5-1:counter-mod}). Nel sistema viene istanziato con \texttt{M} pari a \texttt{N} e \texttt{W} pari a \texttt{A}, con \texttt{load} a '0' e ingresso \texttt{d} a zero; l'uscita \texttt{tc} vale '1' quando il contatore è abilitato e si trova sull'ultimo valore.

L'unità di controllo \texttt{control\_unit} è una FSM di tipo Moore con quattro stati (\texttt{S\_IDLE}, \texttt{S\_READ}, \texttt{S\_WRITE}, \texttt{S\_DONE}), descritta con due process. Il process \texttt{f\_controllo}, sensibile a \texttt{stato\_corrente}, \texttt{start} e \texttt{tc}, assegna dei valori di default alle uscite e calcola lo stato prossimo; il process \texttt{mem\_stato} aggiorna \texttt{stato\_corrente} sul fronte di salita di \texttt{clk}, con reset sincrono verso \texttt{S\_IDLE}. Il comportamento è il seguente:
\begin{enumerate}
  \item in \texttt{S\_IDLE} si attende \texttt{start} a '1' per passare in \texttt{S\_READ};
  \item in \texttt{S\_READ} viene attivato \texttt{rd} e si passa in \texttt{S\_WRITE};
  \item in \texttt{S\_WRITE} vengono attivati \texttt{we} e \texttt{cnt\_en}; se \texttt{tc} vale '1' si va in \texttt{S\_DONE}, altrimenti si torna in \texttt{S\_READ};
  \item in \texttt{S\_DONE} viene attivato \texttt{done}, e si torna in \texttt{S\_IDLE} solo quando \texttt{start} torna a '0'.
\end{enumerate}
Ogni locazione impiega quindi due cicli di clock: nel primo la ROM campiona l'indirizzo, nel secondo il dato elaborato da $M$ è disponibile e viene scritto in MEM nello stesso fronte in cui il contatore avanza. Rispetto alla soluzione proposta, l'unità di controllo ha quattro stati con un altro significato, non ha l'uscita di reset del contatore (che riceve direttamente \texttt{rst}) e ha l'uscita \texttt{done}.

\vhdlfile{esercizio6-1/control_unit.vhd}{Implementazione della Control Unit}{lst:esercizio6-1:cu}

Il top module \texttt{sistema} ha generic \texttt{N} (16) e \texttt{A} (4) e porte \texttt{clk}, \texttt{rst}, \texttt{start} in ingresso, \texttt{done} in uscita, oltre a \texttt{raddr} e \texttt{dout} per la lettura della MEM. L'architettura \texttt{structural} istanzia i cinque componenti con port mapping e dichiara i segnali interni \texttt{addr}, \texttt{rom\_data}, \texttt{m\_data}, \texttt{rd}, \texttt{we}, \texttt{cnt\_en} e \texttt{tc}.

\vhdlfile{esercizio6-1/sistema.vhd}{Implementazione del top module del sistema}{lst:esercizio6-1:sistema}

\simulazione

Il testbench \texttt{sistema\_tb} ha un'entità senza porte e istanzia il top module come \texttt{dut} con \texttt{N} pari a 16 e \texttt{A} pari a 4. Il clock ha periodo di \SI{10}{ns}, generato da un process dedicato. Nel process \texttt{stim} vengono eseguiti questi passi:
\begin{itemize}
  \item \texttt{rst} a '1' per due periodi, poi a '0' per altri due;
  \item \texttt{start} a '1' e attesa di \texttt{done} a '1';
  \item attesa di tre periodi, poi \texttt{start} a '0' e altri due periodi;
  \item lettura di tutte le locazioni della MEM tramite \texttt{raddr}, con un \texttt{assert} per ciascuna, confrontata con i valori attesi della costante \texttt{EXP}, con severità \texttt{error}.
\end{itemize}
I valori attesi sono il numero di '1' di ciascun byte della ROM: $0, 8, 1, 1, 4, 4, 4, 4, 7, 7, 2, 4, 6, 1, 4, 4$.

\vhdlfile{esercizio6-1/sistema_tb.vhd}{Testbench del sistema}{lst:esercizio6-1:tb}

Il risultato della simulazione è riportato in figura~\ref{fig:esercizio6-1:simulazione}.

\begin{figure}[H]
  \centering
  \includegraphics[width=\linewidth]{figure/esercizio6-1/simulazione}
  \caption{Simulazione del sistema PO/PC}
  \label{fig:esercizio6-1:simulazione}
\end{figure}

Dopo il reset iniziale, \texttt{start} sale a '1' intorno a \SI{40}{ns} e la macchina esce da \texttt{S\_IDLE}: \texttt{rd} e \texttt{we} si alternano a ogni ciclo, mentre \texttt{addr} scorre da 0 a 15. Ad esempio, per \texttt{addr} pari a 1 la ROM restituisce 255 (0xFF, $11111111_2$), \texttt{onecount} produce \texttt{m\_data} pari a 8 e la locazione \texttt{ram[1]} passa a 8. Dopo la scrittura dell'ultima locazione lo stato passa a \texttt{S\_DONE} e \texttt{done} sale a '1'; quando \texttt{start} torna a '0' la macchina ritorna in \texttt{S\_IDLE}. Nella fase finale \texttt{raddr} scorre da 0 a 15 e \texttt{dout} riporta il contenuto della MEM: 0, 8, 1, 1, 4, 4, 4, 4, 7, 7, 2, 4, 6, 1, 4, 4, che coincide con \texttt{EXP}.

\begin{risultato}
Il contenuto finale di \texttt{ram} mostrato in figura coincide con i valori attesi di \texttt{EXP}, e \texttt{done} viene attivato una volta al termine della scansione.
\end{risultato}
